\section{Introduction}
Behavioural testing is central to preclinical research\new{. In rodents, it} still relies heavily on short tests performed \emph{outside} the animal's home cage, in a novel environment \cite{kiryk2026}. Such protocols \new{capture} only a narrow window into an animal's activity, are frequently run during the light phase even though rodents are predominantly nocturnal, and require handling and transfer that introduce stress and behavioural bias. \new{This limits the observation of circadian rhythms and slowly emerging changes, and contributes to variability across laboratories} \cite{kiryk2026}.

Home-Cage Monitoring (HCM) inverts this logic: instead of bringing the animal to the test, the observation is brought into the animal's usual environment. Through sensors, hardware, and software, an HCM system continuously observes an animal living in a space that provides food, water, social contact, and safety, over days, weeks or months and with limited human intervention \cite{kiryk2026,huzard2026tech}. This yields more regular, better-contextualised data and supports the 3Rs (Replacement, Reduction, Refinement) --- particularly Refinement --- by reducing handling and handling-related stressors. \new{HCM does not, however, remove methodological bias so much as \emph{shift} part of it toward sensors, analysis software, variable definitions, and metadata quality} \cite{kiryk2026}.

This shift creates a data-interoperability problem. HCM systems vary widely --- video, RFID, infrared and motion sensors, food- and water-intake tracking, or combinations thereof; commercial, lab-built, or hybrid \cite{huzard2026tech,duarte2026diy}. Two studies may measure closely related behaviours while \new{relying on} very different sensors, parameters, and file formats. Crucially, HCM data are inseparable from their experimental context: a measurement is interpretable only if one also knows the animal (species, strain, sex, age, genotype), the enclosure and its environmental conditions, the device and its configuration, the time, and the protocol \cite{forrest2026}. This heterogeneity hinders comparison, reuse, and integration, and stands in the way of FAIR (Findable, Accessible, Interoperable, Reusable) data \cite{fair}. In animal research, data quality is moreover an \emph{ethical} concern, since the data are the durable result of animal experimentation \cite{petit2026wellfair}.

Controlled vocabularies, taxonomies, thesauri, or metadata standards help, but they fall short for a domain whose value lies in the \emph{relations} between an animal, an enclosure, a sensor system, a protocol, environmental conditions, observed behaviours, and provenance. An ontology formalises the concepts of the domain, specifies their relations, reuses existing identifiers, and makes the data queryable, comparable, and integrable, while enabling consistency checking and inference \cite{noy2001}. \new{The HCM community has already structured the domain conceptually: within COST Action TEATIME\footnote{\url{https://www.cost-teatime.org/}} (CA20135), Working Group~2 produced a shared definition of what constitutes an HCM system \cite{teatimeHCMDefinition}, but not a deployable ontology (Sect.~2).}

We present \textbf{HCMO}, the Home-Cage Monitoring Ontology, \new{which formalises HCM as an integrated system and is developed with members of the TEATIME network. Its contributions are:}

\begin{enumerate}
\item \textbf{An ontology for the HCM domain} organised into five modules: a minimal
enclosure-centred \emph{core}, \emph{bio} (subjects, groups, factors, and housing assignments), \emph{env} (environmental profiles and measurements), \emph{obs} (observations and results), and \emph{tech} (hardware, sensors, actuators, software, and time series). The model is built around an explicit \emph{sensor $\neq$ observation $\neq$ result} separation\new{, and keeps the monitored subject (\emph{bio}) and its environmental conditions (\emph{env}) distinct from the observations made about them (\emph{obs})}.

\item \textbf{Selective, reviewed reuse of established standards} \new{(BFO/IAO, SOSA/SSN, QUDT, SemTS, OWL-Time), with an explicit policy on how strongly each reuse is claimed, and an instance-level bridge to the ISA/OBI/STATO ecosystem.}

\item \textbf{A FAIR, tool-consumable resource package}: a stable release manifest,
modular Turtle sources, a reproducibly generated merged graph (TTL/OWL/JSON-LD), SHACL shapes, a JSON-LD context, competency-question SPARQL, HTML documentation, a persistent IRI, an open licence (CC BY 4.0), and a DOI.

\item \textbf{A competency-question-driven evaluation} of coverage and quality\new{, including an executable ISA/RO-Crate round-trip use case}.
\end{enumerate}

\new{The remainder of the paper positions HCMO against related work (Sect.~2), describes the development methodology (Sect.~3), presents the resource (Sect.~4) and its interoperability use case (Sect.~5), reports the evaluation and potential impact (Sect.~6), and discusses limitations and future work (Sect.~7).}
